% Granted invention patents. Define \patentenglish before inputting this file
% to use the English version; otherwise the official Chinese version is used.
\begin{enumerate}[label={[\arabic*]}]
\ifdefined\patentenglish
  \item \textbf{Heng Guo}, Zhanyu Ma, Boxin Shi, Kongming Liang, Yunpeng Jia, Yufei Han. ``3D Reconstruction Method, Apparatus, Device, and Storage Medium.'' Chinese invention patent, ZL202410352394.0, granted October 9, 2026.
  \item \textbf{Heng Guo}, Zhanyu Ma, Ruiyang Chen, Zaiyan Yang. ``A Surface Normal Estimation Method, Apparatus, Device, and Medium for Target Objects.'' Chinese invention patent, ZL202510094535.8, granted June 19, 2026.
  \item \textbf{Heng Guo}, Zhanyu Ma, Mingzhi Pei, Xiangyi Wang. ``A 3D Reconstruction Method, Apparatus, Device, and Medium Based on Multi-view Normal Aggregation.'' Chinese invention patent, ZL202510110808.3, granted April 3, 2026.
  \item \textbf{Heng Guo}, Wenjie Li, Zhanyu Ma. ``Methods for Training an Image Reconstruction Model and Reconstructing Images.'' Chinese invention patent, ZL202510156115.8, granted December 2, 2025.
\else
  \item \textbf{郭亨}，马占宇，施柏鑫，梁孔明，贾云鹏，韩雨霏。《三维重建方法、装置、设备及存储介质》。中国发明专利，专利号：ZL202410352394.0，授权公告日：2026年10月9日。
  \item \textbf{郭亨}，马占宇，陈瑞阳，杨再俨。《一种目标物体的表面法线估计方法、装置、设备及介质》。中国发明专利，专利号：ZL202510094535.8，授权公告日：2026年6月19日。
  \item \textbf{郭亨}，马占宇，裴明志，王象祎。《基于多视角法线聚合的三维重建方法、装置、设备及介质》。中国发明专利，专利号：ZL202510110808.3，授权公告日：2026年4月3日。
  \item \textbf{郭亨}，李文杰，马占宇。《图像重建模型训练以及图像重建方法》。中国发明专利，专利号：ZL202510156115.8，授权公告日：2025年12月2日。
\fi
\end{enumerate}
